\section{Inflación y ondas gravitacionales primordiales en RG}
\label{sec:inflacion}

Esta sección presenta la inflación con un campo escalar y calcula, en RG, el espectro tensorial primordial. Es material estándar \etq{E}, pero se desarrolla con detalle por dos razones: el cálculo de $P_T$ en RG es el punto de referencia contra el que se compara todo lo que sigue, y sus pasos se reutilizan casi sin cambios en gravedad unimodular.

\subsection{Por qué inflación}

El modelo cosmológico estándar sin inflación tiene dos dificultades conocidas. La primera es el problema del horizonte: la radiación cósmica de fondo es homogénea a una parte en $10^5$ en regiones del cielo que, en una expansión decelerada, nunca estuvieron en contacto causal. La segunda es el problema de la planitud: la desviación de la densidad respecto de la crítica satisface $|\Omega-1|\propto 1/(aH)^2$, y si la expansión desacelera, $aH$ decrece, de modo que $|\Omega-1|$ crece con el tiempo; observar hoy un valor pequeño exige un ajuste fino extremo de las condiciones iniciales.

Ambas se resuelven si existió una época en que el \emph{radio de Hubble comóvil} $(aH)^{-1}$ decreció. Como $(aH)^{-1}=1/\dot a$, se cumple
\begin{equation}
\frac{\dd}{\dd t}\,(aH)^{-1}=-\frac{\ddot a}{\dot a^2},
\label{eq:radioHubble}
\end{equation}
que es negativo exactamente cuando $\ddot a>0$: una época de expansión acelerada. Es útil introducir el primer parámetro de flujo de Hubble,
\begin{equation}
\epsilon_1\equiv-\frac{\dot H}{H^2},\qquad\text{de modo que}\qquad \frac{\ddot a}{a}=H^2(1-\epsilon_1),\quad \frac{\dd\ln(aH)}{\dd N}=1-\epsilon_1.
\label{eq:eps1def}
\end{equation}
La inflación es entonces $\epsilon_1<1$, y termina cuando $\epsilon_1$ alcanza $1$. Para que las dos dificultades desaparezcan hace falta que la aceleración dure lo suficiente: del orden de $60$ e-folds \cite{leon22} (con un valor que depende de los detalles del recalentamiento; una estimación estándar sitúa las escalas del fondo cósmico de microondas unos $50$ e-folds antes del final de la inflación \cite{baumann22}).

Hay un segundo aspecto, que es el que hace de la inflación una teoría de estructura y no solo de condiciones iniciales. Un modo de longitud de onda comóvil fija, $\lambda=2\pi/k$, tiene longitud física $a\lambda$, que crece con la expansión, mientras el radio de Hubble físico $H^{-1}$ cambia poco durante la inflación. Un modo que empieza dentro del horizonte ($k\gg aH$) termina fuera de él ($k\ll aH$) al cabo de unos pocos e-folds; el instante $k=aH$ es el \emph{cruce del horizonte}. Las fluctuaciones cuánticas del vacío, que son minúsculas mientras el modo está dentro del horizonte, quedan congeladas al salir y se convierten en perturbaciones clásicas. Ese mecanismo se detalla en la sección~\ref{sec:inf-cuantico}.

\subsection{Inflación con un campo escalar}
\label{sec:inf-campo}

El modelo más simple es un campo escalar $\phi$ (el \emph{inflatón}) con potencial $V(\phi)$ y acción estándar,
\begin{equation}
S_\phi=\int\dd^4x\,\sqrt{-g}\,\Big[-\tfrac12\,g^{ab}\partial_a\phi\,\partial_b\phi-V(\phi)\Big].
\label{eq:Sphi}
\end{equation}
Su tensor energía-impulso, definido por \eqref{eq:Tab}, es $T_{ab}=\partial_a\phi\,\partial_b\phi-g_{ab}\big[\tfrac12(\partial\phi)^2+V\big]$. En el fondo homogéneo, $\phi=\phi(t)$, resulta un fluido perfecto con
\begin{equation}
\rho=\tfrac12\dot\phi^2+V,\qquad p=\tfrac12\dot\phi^2-V,\qquad \rho+p=\dot\phi^2.
\label{eq:rhopphi}
\end{equation}
Las ecuaciones de Friedmann \eqref{eq:Friedmann}, la relación \eqref{eq:Hdot} y la continuidad \eqref{eq:continuidad} se convierten en
\begin{align}
3\MP^2H^2&=\tfrac12\dot\phi^2+V, \label{eq:FriedPhi}\\
\dot H&=-\frac{\dot\phi^2}{2\MP^2}, \label{eq:HdotPhi}\\
\ddot\phi+3H\dot\phi+V'(\phi)&=0. \label{eq:KG}
\end{align}
La \eqref{eq:KG} es la ecuación de Klein--Gordon en un fondo en expansión: el campo rueda por el potencial con una fricción $3H\dot\phi$ que produce la expansión (``fricción de Hubble''). El primer parámetro de flujo es, por \eqref{eq:HdotPhi}, $\epsilon_1=\dot\phi^2/(2\MP^2H^2)$, y la condición de inflación $\epsilon_1<1$ equivale a $\dot\phi^2<V$: la energía potencial domina sobre la cinética.

\paragraph{Slow-roll.} Si el potencial es lo bastante plano, el campo rueda lentamente, $|\ddot\phi|\ll|3H\dot\phi|$, y las ecuaciones se simplifican:
\begin{equation}
\dot\phi\simeq-\frac{V'}{3H},\qquad H^2\simeq\frac{V}{3\MP^2},\qquad \epsilon_1\simeq\epsilon_V\equiv\frac{\MP^2}{2}\Big(\frac{V'}{V}\Big)^2.
\label{eq:slowroll}
\end{equation}
El número de e-folds que quedan hasta el final de la inflación desde un valor $\phi$ es $N=\int H\,\dd t=\int(H/\dot\phi)\,\dd\phi$, es decir, en slow-roll,
\begin{equation}
N(\phi)\simeq\frac{1}{\MP^2}\int_{\phi_e}^{\phi}\frac{V}{V'}\,\dd\phi,
\label{eq:Nslow}
\end{equation}
con $\phi_e$ el valor donde $\epsilon_V=1$.

\paragraph{El ejemplo del potencial cuadrático.} Para $V=\tfrac12 m^2\phi^2$ (el potencial que se usa en el cálculo numérico, sección~\ref{sec:numerico}), $V/V'=\phi/2$ y $\epsilon_V=2\MP^2/\phi^2$, de modo que $\epsilon_V=1$ en $\phi_e^2=2\MP^2$. La integral \eqref{eq:Nslow} da, con $\MP=1$,
\begin{equation}
N\simeq\frac{\phi^2}{4}-\frac12,\qquad \epsilon_V\simeq\frac{1}{2N+1},\qquad H^2\simeq\frac{m^2\phi^2}{6}.
\label{eq:phi2}
\end{equation}
Para $N=60$ resulta $\phi\simeq15.6\,\MP$, $\epsilon_V\simeq0.0083$ y, con $m=6\times10^{-6}\MP$ (el valor que se usa más adelante), $H\simeq3.8\times10^{-5}\MP$. Este potencial, el más simple posible, predice como se verá una relación tensor-escalar $r\simeq16\epsilon_V\simeq0.13$ a $N=60$, en tensión con las cotas actuales; su interés aquí es que es sencillo, analítico y ampliamente conocido, no que sea observacionalmente viable.

\subsection{Fluctuaciones cuánticas: un modelo de juguete}
\label{sec:inf-cuantico}

Antes de las ondas gravitacionales conviene entender el mecanismo en el caso más simple: un campo escalar sin masa $\chi$, minimamente acoplado, sobre un fondo de de Sitter (un universo con $H$ constante, $a=-1/(H\eta)$ con $\eta<0$). Su acción, en tiempo conforme, es $S=\tfrac12\int\dd\eta\,\dd^3x\,a^2[\chi'^2-(\nabla\chi)^2]$. La variable $v=a\chi$ la pone en la forma de un campo con masa dependiente del tiempo,
\begin{equation}
S=\frac12\int\dd\eta\,\dd^3x\,\Big[v'^2-(\nabla v)^2+\frac{a''}{a}\,v^2\Big],
\label{eq:Sv}
\end{equation}
y cada modo de Fourier satisface
\begin{equation}
v_k''+\Big(k^2-\frac{a''}{a}\Big)v_k=0,\qquad\text{con}\qquad \frac{a''}{a}=\frac{2}{\eta^2}\ \text{ en de Sitter.}
\label{eq:vk}
\end{equation}
Es un oscilador armónico con frecuencia dependiente del tiempo, $\omega^2=k^2-a''/a$. Dentro del horizonte, $k^2\gg a''/a$ y el modo oscila con frecuencia $k$, como en el espacio plano; fuera del horizonte, $a''/a$ domina y el modo deja de oscilar.

Se cuantiza el campo promoviendo cada amplitud a un operador, $\hat v=\int\!\frac{\dd^3k}{(2\pi)^3}\big[v_k\,\hat a_{\mathbf k}\,\ee^{i\mathbf k\cdot\mathbf x}+\mathrm{h.c.}\big]$, con la normalización de Wronskiano $v_kv_k^{*\prime}-v_k^*v_k'=i$. Falta elegir el estado. Mucho antes del cruce del horizonte, $|k\eta|\to\infty$, el modo no siente la curvatura y lo natural es el vacío del espacio plano: $v_k\to\ee^{-ik\eta}/\sqrt{2k}$. Es la condición de \emph{Bunch--Davies}. La solución exacta de \eqref{eq:vk} que la cumple es
\begin{equation}
v_k(\eta)=\frac{\ee^{-ik\eta}}{\sqrt{2k}}\Big(1-\frac{i}{k\eta}\Big),
\label{eq:vkBD}
\end{equation}
que se verifica derivando dos veces. El campo original es $\chi_k=v_k/a=-H\eta\,v_k$, y su módulo al cuadrado es
\begin{equation}
|\chi_k|^2=\frac{H^2}{2k^3}\,\big(1+k^2\eta^2\big)\ \longrightarrow\ \frac{H^2}{2k^3}\qquad(|k\eta|\to0).
\label{eq:chi2}
\end{equation}
El espectro de potencia se define como la contribución por intervalo logarítmico de $k$ a la varianza $\langle\chi^2\rangle$; resulta $\mathcal P_\chi=\frac{k^3}{2\pi^2}|\chi_k|^2\to(H/2\pi)^2$. Es el resultado clásico: fuera del horizonte las fluctuaciones de un campo sin masa quedan \emph{congeladas} con una amplitud fijada solo por $H$, y son independientes de $k$ (invariancia de escala) porque en de Sitter $H$ es constante. En un fondo de inflación real, donde $H$ decrece despacio, $H$ se evalúa en el momento del cruce $k=aH$, lo que introduce una débil dependencia con $k$.

\subsection{Ondas gravitacionales primordiales en RG}
\label{sec:GW-RG}

Una onda gravitacional primordial es una perturbación tensorial de la parte espacial de la métrica,
\begin{equation}
g_{ij}=a^2\,(\delta_{ij}+h_{ij}),\qquad h_{ii}=0,\qquad \partial_ih_{ij}=0,
\label{eq:hij}
\end{equation}
sin traza y transversa (TT). Una matriz simétrica $3\times3$ tiene seis componentes; anular la traza quita una y la transversalidad quita tres, de modo que quedan \textbf{dos polarizaciones}, $+$ y $\times$. La acción cuadrática en $h_{ij}$ (que se deriva con detalle en la sección~\ref{sec:accion2}, y que en RG es un resultado estándar) es
\begin{equation}
S^{(2)}=\frac{\MP^2}{8}\int\dd t\,\dd^3x\,\Big[a^3\,\dot h_{ij}\dot h_{ij}-a\,(\partial_kh_{ij})^2\Big].
\label{eq:S2RG}
\end{equation}
Se descompone en modos de Fourier y polarizaciones, $h_{ij}=\sum_{\lambda=+,\times}\int\!\frac{\dd^3k}{(2\pi)^3}\,h_\lambda(\eta,\mathbf k)\,e^\lambda_{ij}(\mathbf k)\,\ee^{i\mathbf k\cdot\mathbf x}$, con tensores de polarización normalizados por $e^\lambda_{ij}e^{\lambda'*}_{ij}=2\delta_{\lambda\lambda'}$. Esa normalización es una convención, y la elección se discute en la sección~\ref{sec:normalizacion}. Con ella, la acción de cada polarización, en tiempo conforme, es $\frac{\MP^2}{4}\int\dd\eta\,a^2\,[|h_\lambda'|^2-k^2|h_\lambda|^2]$, y la variable canónica
\begin{equation}
v_\lambda\equiv\frac{a\,\MP}{\sqrt2}\,h_\lambda
\label{eq:vcanon}
\end{equation}
la deja en la forma $\tfrac12\int\dd\eta\,[|v_\lambda'|^2-(k^2-a''/a)|v_\lambda|^2]$: exactamente la del campo escalar sin masa de la sección anterior. Cada polarización se comporta, por lo tanto, como un campo escalar sin masa, y todo lo dicho se aplica: obedece $v_\lambda''+(k^2-a''/a)v_\lambda=0$, se cuantiza con la condición de Bunch--Davies, y en de Sitter su módulo cuadrado es el de \eqref{eq:chi2}, es decir $|h_\lambda|^2=2|v_\lambda|^2/(a^2\MP^2)\to H^2/(k^3\MP^2)$ fuera del horizonte.

El espectro tensorial se define como la contribución por intervalo logarítmico a la varianza de la perturbación completa, $\langle h_{ij}h_{ij}\rangle=\int\dd\ln k\,P_T(k)$. Como $e^\lambda_{ij}e^{\lambda*}_{ij}=2$, resulta $P_T=\frac{k^3}{2\pi^2}\cdot2\sum_\lambda|h_\lambda|^2$, y con el resultado de de Sitter (dos polarizaciones iguales),
\begin{equation}
\boxed{\;P_T(k)=\frac{2H^2}{\pi^2\MP^2}\Big|_{k=aH}\;}
\label{eq:PTRG}
\end{equation}
Esta es la fórmula estándar \etq{E}, y es la convención con la que se define $r$ \cite{baumann22,martin14}. Para un fondo de inflación de slow-roll, $H$ varía lentamente, de modo que $P_T$ tiene una pendiente
\begin{equation}
n_T\equiv\frac{\dd\ln P_T}{\dd\ln k}\simeq-2\epsilon_1,
\label{eq:nT}
\end{equation}
ligeramente roja, porque los modos de mayor $k$ salen del horizonte más tarde, cuando $H$ es menor. Usando $H^2\simeq V/(3\MP^2)$, la amplitud puede escribirse $P_T\simeq2V/(3\pi^2\MP^4)$, que muestra por qué medir el espectro tensorial es medir la escala de energía de la inflación \cite{baumann22}. Los primeros términos de la corrección de slow-roll son \cite{martin14}
\begin{equation}
P_T=\frac{2H^2}{\pi^2\MP^2}\Big[1-2(C+1)\,\epsilon_1+\mathcal O(\epsilon^2)\Big]_{k=aH},\qquad C\equiv\gamma_{\rm E}+\ln2-2\simeq-0.7296,
\label{eq:PTNLO}
\end{equation}
con $\gamma_{\rm E}$ la constante de Euler, y las correcciones a segundo orden se usan en los controles numéricos de la sección~\ref{sec:resultados}.

Falta la otra mitad de $r$. El espectro de las perturbaciones escalares de curvatura, cuya derivación es más larga y no se reproduce aquí, es \cite{martin14}
\begin{equation}
P_\mathcal{R}=\frac{H^2}{8\pi^2\epsilon_1\MP^2}\Big|_{k=aH},\qquad r\equiv\frac{P_T}{P_\mathcal{R}}\simeq16\,\epsilon_1.
\label{eq:PRr}
\end{equation}
Con los datos actuales (el análisis de Planck 2018 con BICEP/Keck 2014 da $r<0.064$ al 95\% \cite{planck20}, y BICEP/Keck hasta la temporada de observación de 2018 da $r<0.036$ \cite{bicep21}, valor que solo se leyó en el resumen del artículo), un potencial puramente cuadrático, que da $r\approx0.13$--$0.16$ para $N_*=60$--$50$, está en tensión con las observaciones: Planck lo describe como un caso que ``siempre predice una relación tensor-escalar grande (de aproximadamente $0.15$), en tensión con los datos'' (traducción libre de \cite{planck20}). Que sea observacionalmente disfavorecido no afecta a su uso en este trabajo, cuyo objetivo es comparar RG y UG con el mismo potencial.
